\SourceLeaf{225}%
der halbe Umfang eines Kreises, dessen Halbmesser $=r$:

\[
=\tfrac{1}{2}\pi k\bigl(e^{\frac{r}{k}}-e^{-\frac{r}{k}}\bigr),
\]

\noindent wo $k$ eine Constante ist, von der wir durch Erfahrung wissen, dass sie gegen alles durch uns messbare ungeheuer gross sein muss. In Euklids Geometrie wird sie unendlich.

In der \textls{Bildersprache} des Unendlichen würde man also sagen müssen, dass die Peripherien zweier unendlichen Kreise, deren Halbmesser um eine endliche Grösse verschieden sind, selbst um eine Grösse verschieden sind, die zu ihnen ein endliches Verhältniss hat.

Hierin ist aber nichts widersprechendes, wenn der endliche Mensch sich nicht vermisst, etwas Unendliches als etwas Gegebenes und von ihm mit seiner gewohnten Anschauung zu Umspannendes betrachten zu wollen.

Sie sehen, dass hier in der That der Fragepunkt unmittelbar an die Metaphysik streift.

\bigskip
\rule{4cm}{0.4pt}
\bigskip

\Needspace{7\baselineskip}
\begin{center}
[\textsc{Schumacher} \emph{an} \textsc{Gauss}.\ \emph{Altona, 19. Juli 1831.}]
\end{center}

\smallskip
\{ Meinen herzlichsten Dank statte ich Ihnen, mein theuerster Freund, für Ihren letzten Brief ab. Ich kann nicht sagen, dass er mich schon überzeugt hätte. Ich glaube die unendliche Grösse nicht als geschlossen gebraucht zu haben. Mir scheint, man kann zeigen, dass mit dem Wachsen des Halbmessers die Differenz der Winkelpunkte des Dreiecks immer mehr verschwindet, und sich der Grenze des Zusammenfallens, so viel man immer will, nähert. Sagt man also, der Kürze halber, sie fallen für einen unendlichen Radius wirklich zusammen, so wird diess eben so wie gewöhnlich verstanden, und es folgt daraus, dass, in Bezug auf die Peripherie, die von den geraden Linien intercaptir%
